\documentclass[10pt,a4paper]{article}
\usepackage{amsmath,amssymb,amsthm,graphicx}
\usepackage{xfrac}
\newenvironment{enumeratei}
{\bgroup\def\theenumi{\roman{enumi}}\def\theenumii{\arabic{enumii}}\begin{enumerate}}
{\end{enumerate}\egroup}

\hyphenation{addi-tion-al belong belongs denote exchan-ge exchan-ged exchan-ges}

\newcommand{\psfrac}[2]{\sfrac{(#1)}{#2}}

\newcommand\mto{\mathchoice{\longmapsto}{\mapsto}{\mapsto}{\mapsto}}
\let\oldbigsqcup\bigsqcup
\renewcommand{\bigsqcup}{\mathop{\textstyle\oldbigsqcup}\displaylimits}

\let\epsilon\varepsilon

\DeclareMathOperator{\Ess}{Ess}
\DeclareMathOperator{\Leb}{Leb}

\usepackage{tikz}
\usetikzlibrary{plotmarks}

\theoremstyle{plain}
\newtheorem{theorem}{Theorem}[section]
\newtheorem{lemma}[theorem]{Lemma}
\newtheorem{proposition}[theorem]{Proposition}
\newtheorem{corollary}[theorem]{Corollary}
\newtheorem{claim}{Claim}

\theoremstyle{definition}
\newtheorem{remark}[theorem]{Remark}
\newtheorem{example}[theorem]{Example}

\newcommand{\A}{\mathcal{A}}
\newcommand{\B}{\mathcal{B}}
\newcommand{\I}{\mathcal{I}}

\newcommand{\Z}{\mathbb{Z}}
\newcommand{\N}{\mathbb{N}}
\newcommand{\R}{\mathbb{R}}
\title{The induced transformation on a symmetric subinterval satisfies the same conjugacy equation}
\author{Przemyslaw Berk, Frank Trujillo}
\date{doi:10.5802/jep.302; CC BY4.0}
\begin{document}
\maketitle
\noindent\textit{Editorial note.} The counted statement is the author's labelled lemma environment at source0.tex lines 626-628; its proof environment opens at line 630, the line after the blank separator, so the pairing is by adjacency and the retained proof range 630-646 is the whole environment.\par\medskip
\noindent\textit{Editorial note.} Two statements are retained as context because the counted statement and proof refer to them: the conjugacy equation eq: standardconjugacy (source lines 463-465), which the statement itself displays, and the local reflection lemma (source lines 596-602), which the proof cites at source line 639.\par\medskip
\noindent\textit{Editorial note.} This source is new to the material-authoring work and its preamble is built from the author's own macro and theorem block (source lines 8-45) rather than from a prepared adapter.\par\medskip
\noindent\textit{Editorial note.} This source has no label-number dump and the retained counter values are left empty, so the printed numbering of the card is the wrapper's default; the identity of the counted item is its author label.\par\medskip
\noindent\textit{Editorial note.} r002 note: r001 failed with an undefined reference to eq:generalized\_conjugacy, which the counted proof applies at source lines 631 and 641; r002 adds exactly that numbered identity as a further context and changes nothing else.\par\medskip
\noindent\textit{Editorial scope.} One original result of the pinned publisher article, counted once: the induced transformation on a symmetric subinterval satisfies the conjugacy equation (author label lem:symmetric\_condition).\par\medskip
\noindent\textit{Allowed standard prerequisites (editorial).}\begin{itemize}
\item Induced transformation. The induced transformation on a subinterval is defined by the first return time of its points to the subinterval.. It is the transformation about which the statement is made.
\item Conjugacy equation. For a symmetric transformation the reflection of the whole interval conjugates it to its inverse off the discontinuities.. It is the equation the induced transformation is asserted to satisfy.
\item Symmetry of subintervals. A subinterval is either alpha-symmetric for a continuity interval or symmetric about the middle of the whole interval.. It is the case distinction along which the proof runs.
\end{itemize}\medskip
\begin{multline}
\label{eq:generalized_conjugacy}
\mathcal{I}_I\circ T^n(x)=T^{-n}\circ\mathcal{I}_I(x) \quad \text{if } x \neq T^{-i}(\partial I_\alpha) \text{ for any } \alpha \in \A \text{ and}\\ \text{any $i$ such that } 0 \leq i < n.
\end{multline}

\begin{equation}\label{eq: standardconjugacy}
\mathcal{I}_I\circ T(x)=T^{-1}\circ\mathcal{I}_I(x) \quad \text{if } x \neq \partial I_{\alpha} \text{ for any } \alpha \in \A.
\end{equation}

\begin{lemma}\label{lem: local_reflection}
Let $T : I \to I$ be a symmetric IET and $J \subseteq I$ be an $\alpha$-symmetric interval, for some $\alpha \in \A$. Then $\mathcal{I}_I\circ T(\mathring{J}) = \mathring{J}$. Moreover,
\[
\I_J = \mathcal{I}_I\circ T|_{\mathring{J}},
\]
that is, $\mathcal{I}_I\circ T|_{\mathring{J}}$ is the symmetric reflection on $J$.
\end{lemma}

\begin{lemma}
\label{lem:symmetric_condition}
Let $T : I \to I$ be a symmetric IET and $J \subseteq I$ be a symmetric interval. Then $T_J$ satisfies \eqref{eq: standardconjugacy}.\end{lemma}

\begin{proof}
Assume first that $J$ is $\alpha$-symmetric for some $\alpha\in\A$. Let $x\in J$ be not an endpoint of an interval exchanged by $T_J$ and let $h\in\N$ be such that $T_J(x)=T^{h}(x)$. Notice that since $x$ is not an endpoint of the exchanged intervals, by~\eqref{eq:generalized_conjugacy} applied to~$T$,\vspace*{-3pt}
\[
{(\I_I\circ T)\circ T_{J}(x)} = (\I_I\circ T) \circ T^h(x)=T^{-h}\circ (\I_I\circ T)(x).
\]
Thus, to prove \eqref{eq: standardconjugacy}, it~suffices to notice that\vspace*{-3pt}
\[
T_{J}^{-1}\circ (\I_I\circ T)(x) = T^{-h}\circ (\I_I\circ T)(x).
\]
Indeed, by~Lemma \ref{lem: local_reflection}, for any $n \geq 1$, $T^n(x) \in J$ if and only if $(\I_I\circ T) \circ T^n(x) \in J$. Thus the equation above follows since $T_J(x)=T^{h}(x)$ and $T^{-h}\circ (\I_I\circ T)(x) = (\I_I\circ T) \circ T^h(x)$.

Assume now that $J$ is $\sfrac{1}{2}$-symmetric. Let $x\in J$ be not an endpoint of an interval exchanged by $T_J$ and let $h\in\N$ be such that $T_J(x)=T^{h}(x)$. Again by \eqref{eq:generalized_conjugacy} we get\vspace*{-3pt}
\[
\I_I\circ T_{J}(x) = \I_I\circ T^h(x)=T^{-h}\circ \I_I(x)
\]
and, similarly to the previous case, to show \eqref{eq: standardconjugacy} it is sufficient to notice that $T_{J}^{-1}\circ\nobreak\I_I(x)=T^{-h}\circ\I_I(x)$.
\end{proof}

\end{document}
